\documentclass[12pt]{article}
\usepackage[utf8]{inputenc}
\usepackage[spanish]{babel}
\usepackage{amsmath}
\usepackage{amssymb}
\usepackage{geometry}
\usepackage{graphicx}
\usepackage{xcolor}
\geometry{a4paper, margin=1in}

\title{\textbf{Notas de Clase: Ondas Electromagnéticas y Óptica Coherente}}
\author{}
\date{}

\begin{document}
	
	\maketitle
	
	\section{Características de las Ondas Electromagnéticas (EM)}
	\begin{enumerate}
		\item Cuando una carga eléctrica vibra, el campo eléctrico a su alrededor cambia, creando a su vez un campo magnético cambiante.
		\item Los campos magnético y eléctrico se inducen mutuamente de forma continua.
		\item Una onda EM viaja en todas direcciones en el espacio, es decir, son tridimensionales.
		\item Los campos eléctricos y magnéticos vibran en ángulo recto (perpendicularmente) con respecto a la dirección en la que viaja la onda, por lo que se define como una onda transversal.
		\item Toda la materia contiene partículas cargadas en constante movimiento; por lo tanto, todos los objetos emiten ondas EM.
		\item Las longitudes de onda emitidas se acortan a medida que aumenta la temperatura del material emisor.
		\item Las ondas EM transportan energía radiante a través del espacio.
		\item La luz no se produce como una onda continua ininterrumpida, sino que está formada por cuantos de energía llamados fotones, que se emiten continuamente en forma de haces.
	\end{enumerate}
	
	\subsection{Representación en Notación Compleja}
	\textit{Las ondas electromagnéticas son ondas transversales compuestas por campos eléctricos y magnéticos entrelazados que oscilan, cuyas fases se miden en unidades angulares (radianes o grados).}
	
	\vspace{0.5cm}
	\begin{minipage}{0.45\textwidth}
		\begin{align*}
			E &= E_0 e^{i[\varphi(z) \pm (\omega t)]} \\
			E_0 &= \sqrt{\frac{2P_0}{\varepsilon c A}} \\
			\varphi(z) &= \frac{2\pi z n}{\lambda} + \theta_0 = \kappa z n + \theta_0 \\
			\omega &= 2\pi f \\
			\lambda &= \frac{c}{f} = c T \\
			\kappa &= \frac{2\pi}{\lambda}
		\end{align*}
	\end{minipage}
	\hfill
	\begin{minipage}{0.5\textwidth}
		\textbf{Definición de variables:}
		\begin{itemize}
			\item $E_0$: Amplitud de la onda.
			\item $P_0$: Potencia óptica del haz.
			\item $\varepsilon$: Permitividad del medio.
			\item $A$: Área transversal del haz.
			\item $c$: Velocidad de la luz en el vacío ($3\times10^8$ m/s).
			\item $\omega$: Frecuencia angular (rad/s).
			\item $\lambda$: Longitud de onda (distancia entre crestas consecutivas).
			\item $f$: Frecuencia (ciclos por segundo o Hz). Si $f$ incrementa, $\lambda$ disminuye.
			\item $\varphi(z)$: Fase instantánea espacial.
			\item $z$: Longitud del camino físico recorrido.
			\item $\theta_0$: Fase inicial constante de la onda.
			\item $n$: Índice de refracción del medio.
			\item $\kappa$: Número de onda (rad/m).
		\end{itemize}
	\end{minipage}
	
	\subsection{Fase Instantánea y Amplitud}
	\textbf{La fase instantánea espacial, $\varphi(z)$, es proporcional a la distancia recorrida a lo largo del camino óptico.} El camino óptico es la trayectoria que sigue la luz entre dos puntos, definido como el producto de la distancia geométrica real ($z$) y el índice de refracción del medio ($n$).
	
	Partiendo de la fórmula de Euler, la ecuación de onda se expande como:
	\begin{equation*}
		E = E_0 e^{i[\varphi(z) \pm \omega t]} = E_0 e^{i\theta} = E_0 (\cos\theta + i\sin\theta)
	\end{equation*}
	
	Donde definimos la fase total $\theta$:
	\begin{align*}
		\theta &= \varphi(z) \pm \omega t = \kappa z n + \theta_0 \pm \omega t
	\end{align*}
	Aquí, $\omega t$ es la \textbf{dependencia temporal} y $\varphi(z) = \kappa z n + \theta_0$ es la \textbf{dependencia espacial}.
	
	Si expresamos el campo en sus componentes real e imaginaria ($E = x + iy$):
	\begin{align*}
		x &= E_0 \cos\theta \\
		y &= E_0 \sin\theta
	\end{align*}
	
	\textbf{Recuperación de la fase:}
	La información de la fase se recupera relacionando las componentes:
	\begin{equation*}
		\tan\theta = \frac{y}{x} = \frac{\text{Parte Imaginaria}}{\text{Parte Real}} = \frac{\sin\theta}{\cos\theta} \implies \theta = \arctan\left(\frac{y}{x}\right)
	\end{equation*}
	
	\textbf{Recuperación de la amplitud:}
	La amplitud física al cuadrado se obtiene multiplicando el campo por su complejo conjugado ($E^*$):
	\begin{equation*}
		E_0^2 = E E^* = (x+iy)(x-iy) = x^2 + y^2 \implies E_0 = \sqrt{x^2 + y^2}
	\end{equation*}
	
	Físicamente, los campos son reales. Si tomamos la parte real de la ecuación de onda:
	\begin{equation*}
		E = Re(E_0 e^{i\theta}) = E_0 \cos(\varphi(z) \pm \omega t)
	\end{equation*}
	\textit{Nota:} También es válido representarlo con la función seno, dado que un desfase de $\pi/2$ las hace equivalentes: $E = E_0 \sin(\theta + \pi/2) = E_0 \cos(\theta)$.
	
	\subsection{Representación Gráfica}
	Fijando el tiempo ($t=0$), observamos la onda en el espacio:
	\begin{equation*}
		E(x,0) = E_0 \sin\left(\frac{2\pi x}{\lambda}\right)
	\end{equation*}
	\textit{[Imagen: diagrama\_onda\_senoidal\_espacio.png]}
	
	Fijando la posición ($x=0$), observamos la oscilación en el tiempo:
	\begin{equation*}
		E(0,t) = E_0 \sin(\omega t)
	\end{equation*}
	\textit{[Imagen: diagrama\_oscilacion\_tiempo.png]}
	
	\section{Naturaleza de la Luz y Propagación}
	
	\subsection{Dualidad Onda-Partícula}
	Desde la perspectiva de la física moderna, la luz presenta un comportamiento dual:
	\begin{itemize}
		\item \textbf{Naturaleza ondulatoria:} La luz consiste en oscilaciones de campos eléctricos y magnéticos perpendiculares entre sí, capaces de propagarse en el vacío. Se manifiesta en fenómenos como interferencia y difracción.
		\item \textbf{Naturaleza corpuscular:} Se comporta como un flujo de partículas discretas o "cuantos" de energía llamados fotones.
		\item \textbf{Visión heurística:} Cuando el campo electromagnético es de muy baja densidad, dominan las interacciones individuales (fotón-átomo), requiriendo un análisis discreto/cuántico. En campos de alta densidad, con un número enorme de fotones, el fenómeno promedia sus efectos manifestándose como un continuo, permitiendo un análisis ondulatorio clásico. (Similar al agua: moléculas discretas a nivel microscópico, fluido continuo a nivel macroscópico).
	\end{itemize}
	
	\subsection{Ecuación de Onda 3D y Ondas Planas}
	Para una onda que se propaga en cualquier dirección del espacio tridimensional:
	\begin{equation*}
		E(x,y,z,t) = E_0 e^{i[\vec{k}\cdot\vec{r} \pm \omega t + \theta_0]}
	\end{equation*}
	
	\begin{center}
		\fbox{
			\begin{minipage}{0.85\textwidth}
				\textbf{Parámetros Vectoriales:}
				\begin{align*}
					\vec{k} &= (k_x, k_y, k_z) \quad \text{(Vector de onda, indica dirección de propagación)} \\
					\vec{r} &= (x, y, z) \quad \text{(Vector de posición en el espacio)}
				\end{align*}
				\begin{equation*}
					|\vec{k}| = \sqrt{k_x^2 + k_y^2 + k_z^2} = \frac{2\pi}{\lambda}
				\end{equation*}
				\begin{equation*}
					\vec{k}\cdot\vec{r} = k_x x + k_y y + k_z z
				\end{equation*}
			\end{minipage}
		}
	\end{center}
	
	\vspace{0.5cm}
	\textbf{El concepto de "Onda Plana"}\\
	Una onda plana es un modelo idealizado donde la magnitud física (como la fase) es constante sobre cualquier plano perpendicular a la dirección de propagación. La condición geométrica para estos planos es:
	\begin{equation*}
		\vec{k}\cdot\vec{r} = \text{constante}
	\end{equation*}
	A estos contornos de fase constante se les llama \textbf{frentes de onda}.
	
	\begin{itemize}
		\item Se desplazan a la velocidad de la luz en la dirección del vector $\vec{k}$.
		\item Están separados uniformemente por una distancia igual a una longitud de onda ($\lambda$).
		\item El vector campo eléctrico $\vec{E}_0$, el campo magnético $\vec{H}_0$, y el vector $\vec{k}$ son mutuamente perpendiculares.
	\end{itemize}
	
	\textit{[Imagen: diagramas\_frentes\_de\_onda.png]}
	
	\subsubsection*{Rayos láser vs. Ondas planas ideales}
	Una onda plana perfecta posee frentes de onda infinitos y, por ende, energía infinita; no existe en la realidad física.
	
	\textit{[Imagen: frentes\_onda\_planos\_espacio.png]}
	
	En contraste, un rayo láser está espacialmente confinado. Matemáticamente, se aproxima multiplicando la propagación de una onda plana en la dirección $z$ por una distribución gaussiana de intensidad en el plano transversal ($x,y$):
	\begin{equation*}
		E(x,y,z,t) = E_0 \exp\left[-\frac{x^2+y^2}{w^2}\right] \exp[i(kz - \omega t)]
	\end{equation*}
	Donde $w$ es el radio del haz (waist).
	
	\textit{[Imagen: frentes\_onda\_localizadas\_haz\_laser.png]}
	
	\section{Superposición e Interferencia de Ondas}
	
	\subsection{Principio de Superposición}
	Si varias ondas viajan simultáneamente a través de un medio, la perturbación resultante en cualquier punto es la suma vectorial (algebraica) de las perturbaciones individuales de cada onda en ese instante.
	
	\textbf{Relación de Fase:}
	El resultado no solo depende de las amplitudes, sino del desfase relativo:
	\begin{itemize}
		\item \textbf{En fase:} Las crestas coinciden, reforzando la amplitud (suma).
		\item \textbf{Desfasadas ($180^\circ$ o $\pi$ rad):} La cresta de una coincide con el valle de la otra, resultando en atenuación o cancelación (resta).
	\end{itemize}
	\textit{[Imagen: sumas\_y\_cancelacion\_amplitudes.png]}
	
	\subsubsection*{Superposición de Dos Ondas Armónicas (Ondas Estacionarias)}
	Si superponemos dos ondas de igual amplitud y frecuencia que viajan en direcciones opuestas:
	\begin{align*}
		E_1 &= E\sin(kx - \omega t) \quad \text{(Viaja a la derecha)} \\
		E_2 &= E\sin(kx + \omega t) \quad \text{(Viaja a la izquierda)}
	\end{align*}
	El campo resultante es $E_R = E_1 + E_2$. Utilizando la identidad trigonométrica $\sin(A\pm B) = \sin A \cos B \pm \cos A \sin B$:
	\begin{equation*}
		E_R = 2E\sin(kx)\cos(\omega t)
	\end{equation*}
	Esto describe una onda estacionaria, donde la amplitud espacial máxima es $2E$ en los puntos donde $\sin(kx) = 1$ (vientres).
	
	\subsubsection*{Superposición con Diferente Amplitud y Desfase General}
	Dadas dos ondas de misma frecuencia propagándose hacia el mismo punto:
	\begin{align*}
		\vec{E}_1 &= E_{01}\cos(\omega t + kr_1)\hat{r} \\
		\vec{E}_2 &= E_{02}\cos(\omega t + kr_2)\hat{r}
	\end{align*}
	
	El \textbf{desfase inicial} entre ambas debido a la diferencia de recorrido es:
	\begin{equation*}
		\Delta\theta = k(r_1 - r_2) = \frac{2\pi}{\lambda}(r_1 - r_2)
	\end{equation*}
	
	\begin{itemize}
		\item \textbf{En fase:} $\Delta\theta = 2m\pi$ (donde $m$ es entero).
		\item \textbf{Desfasadas (Antifase):} $\Delta\theta = (2m+1)\pi$. \textit{(Nota de corrección respecto a tus notas originales donde decía $\pi/2$; el desfase destructivo total ocurre a múltiplos impares de $\pi$).}
	\end{itemize}
	\textit{[Imagen: representacion\_in\_phase\_out\_of\_phase.png]}
	
	La amplitud resultante $E_0$ se obtiene de la suma vectorial (o empleando notación compleja $E_0^2 = E_T E_T^*$):
	\begin{align*}
		E_0^2 &= E_{01}^2 + E_{02}^2 + 2E_{01}E_{02}\cos(\Delta\theta) \\
		E_0 &= \sqrt{E_{01}^2 + E_{02}^2 + 2E_{01}E_{02}\cos(\Delta\theta)}
	\end{align*}
	De modo que la onda final se puede escribir como $\vec{E} = E_0 \cos(\omega t - \theta)\hat{r}$.
	
	\subsection{Interferencia}
	Se denomina interferencia a la superposición coherente (espacial y temporal) de dos o más frentes de onda en un punto, generando un \textbf{patrón de franjas de intensidad} (máximos y mínimos). 
	
	Existen dos métodos principales para producir luz capaz de interferir a partir de una fuente única:
	\begin{enumerate}
		\item \textbf{División de amplitud:} Un divisor de haz (espejo semitransparente) refleja parte de la luz y transmite el resto (ej. Interferómetro de Michelson).
		\item \textbf{División de frente de onda:} Un solo frente pasa a través de aperturas separadas espacialmente creando frentes secundarios (ej. Experimento de la doble rendija de Young).
	\end{enumerate}
	
	\textit{[Imagen: patron\_interferencia\_dos\_fuentes.png]}
	\textit{[Imagen: graficas\_interferencia\_constructiva\_destructiva\_normal.png]}
	
	\subsubsection*{Condiciones de Interferencia}
	\begin{center}
		\begin{tabular}{|p{0.45\textwidth}|p{0.45\textwidth}|}
			\hline
			\textbf{Interferencia Constructiva} & \textbf{Interferencia Destructiva} \\
			\hline
			Se refuerza el movimiento ondulatorio. Ocurre cuando $\cos(\Delta\theta) = 1$. & Se atenúa o cancela la onda. Ocurre cuando $\cos(\Delta\theta) = -1$. \\
			Diferencia de fase: $\Delta\theta = 2\pi m$ & Diferencia de fase: $\Delta\theta = (2m+1)\pi$ \\
			Diferencia de camino óptico: $\text{OPD} = m\lambda$ & Diferencia de camino óptico: $\text{OPD} = (m + 0.5)\lambda$ \\
			\hline
		\end{tabular}
	\end{center}
	\textit{[Imagen: diagramas\_interferencia\_dos\_fuentes\_constructiva\_destructiva.png]}
	
	\textbf{Propiedad General:} Para generar un patrón de interferencia estable, las fuentes deben ser \textbf{coherentes}; es decir, deben mantener una diferencia de fase constante en el tiempo (y poseer la misma frecuencia).
	
	\section{Intensidad (I) de la Radiación Electromagnética}
	El campo eléctrico de la luz oscila a frecuencias altísimas ($\sim 10^{14}$ Hz). Ningún sensor físico (ojo, película, fotodiodos) responde lo suficientemente rápido para medir la amplitud instantánea. Lo que se registra es la **Irradiancia** (flujo de energía temporal promedio por unidad de área), la cual es proporcional al cuadrado de la amplitud.
	
	\begin{equation*}
		I \propto \langle |E_T|^2 \rangle \implies I = E_T E_T^*
	\end{equation*}
	
	Dimensiones de Intensidad ($W/m^2$):
	\begin{equation*}
		I = \frac{\text{Potencia}}{\text{Área}} = \left(\frac{\text{Energía}}{\text{Volumen}}\right) \times (\text{Velocidad de la onda})
	\end{equation*}
	
	Integrando la energía cinética máxima de los osciladores armónicos en el medio electromagnético:
	\begin{equation*}
		I = \frac{1}{2} \rho c (E_0 \omega)^2
	\end{equation*}
	\textit{[Imagen: grafica\_dos\_ondas\_amplitudes\_diferentes.png]}
	
	\subsection{Interferencia e Intensidad (Ecuación Fundamental)}
	Cuando dos ondas complejas $E_1 = E_{01}e^{i\varphi_1}$ y $E_2 = E_{02}e^{i\varphi_2}$ se superponen, la intensidad observada no es simplemente la suma de las intensidades, debido al término cruzado de interferencia.
	
	\begin{align*}
		I &= (E_1 + E_2)(E_1^* + E_2^*) \\
		&= E_1E_1^* + E_2E_2^* + E_1E_2^* + E_2E_1^* \\
		&= I_1 + I_2 + E_{01}E_{02}e^{i(\varphi_1 - \varphi_2)} + E_{01}E_{02}e^{-i(\varphi_1 - \varphi_2)}
	\end{align*}
	
	Utilizando la identidad de Euler $e^{ix} + e^{-ix} = 2\cos(x)$, llegamos a la Ley de Interferencia:
	\begin{center}
		\fbox{\Large $\mathbf{I = I_1 + I_2 + 2\sqrt{I_1 I_2} \cos(\Delta\theta)}$}
	\end{center}
	Donde $\Delta\theta = \varphi_1 - \varphi_2$.
	
	El término $2\sqrt{I_1 I_2}\cos(\Delta\theta)$ modula la intensidad generando el patrón de franjas, fluctuando entre:
	\begin{itemize}
		\item $I_{max} = I_1 + I_2 + 2\sqrt{I_1 I_2} = (\sqrt{I_1} + \sqrt{I_2})^2$
		\item $I_{min} = I_1 + I_2 - 2\sqrt{I_1 I_2} = (\sqrt{I_1} - \sqrt{I_2})^2$
	\end{itemize}
	\textit{[Imagen: grafica\_intensidad\_patron\_franjas.png]}
	
	Si $I_1 = I_2 = I_0$, la modulación es del 100\%:
	\begin{equation*}
		I = 2I_0 [1 + \cos(\Delta\theta)] = 4I_0 \cos^2\left(\frac{\Delta\theta}{2}\right)
	\end{equation*}
	
	\section{Camino Geométrico y Camino Óptico}
	
	\begin{itemize}
		\item \textbf{Camino Geométrico ($z$):} Distancia física real que recorre la luz medida en metros.
		\item \textbf{Camino Óptico ($L$):} Distancia equivalente que recorrería la luz en el vacío durante el mismo intervalo de tiempo. Se calcula como: $L = n \times z$.
		\item \textbf{Diferencia de Camino Óptico (OPD):} Es la resta entre los caminos ópticos de dos haces interferentes: $\text{OPD} = L_2 - L_1 = n_2z_2 - n_1z_1$.
	\end{itemize}
	
	La diferencia de fase se relaciona directamente con el OPD:
	\begin{equation*}
		\Delta\theta = \frac{2\pi}{\lambda_0} \times \text{OPD} = k \times \text{OPD}
	\end{equation*}
	\textit{[Imagen: esquema\_caminos\_optico\_y\_geometrico\_punto\_Q.png]}
	
	\subsection{Ejemplo de Medios con Diferente Índice de Refracción}
	Dos haces coherentes viajan una misma distancia geométrica $z$. El haz 1 viaja por aire ($n_1 \approx 1$) y el haz 2 por acrílico ($n_2 = 1.5$).
	\textit{[Imagen: esquema\_dos\_haces\_aire\_acrilico\_punto\_Q.png]}
	
	Aunque la diferencia de camino geométrico es $z - z = 0$, la velocidad disminuye dentro del acrílico, acortando su longitud de onda efectiva ($\lambda_{ac} = \lambda_{air}/1.5$).
	\textit{[Imagen: esquema\_ondas\_acrilico\_reduccion\_velocidad.png]}
	
	La diferencia de camino óptico que genera el desfase es:
	\begin{equation*}
		\text{OPD} = n_2 z - n_1 z = z(1.5 - 1) = 0.5z
	\end{equation*}
	
	\subsection{Visibilidad o Contraste de Franjas ($\gamma$)}
	Propuesta por Michelson, cuantifica la calidad y nitidez del patrón de interferencia. Es una variable adimensional entre 0 y 1.
	
	\begin{equation*}
		\gamma = \frac{I_{max} - I_{min}}{I_{max} + I_{min}}
	\end{equation*}
	
	Sustituyendo las intensidades óptimas, la visibilidad se puede expresar como:
	\begin{equation*}
		\gamma = \frac{2\sqrt{I_1 I_2}}{I_1 + I_2}
	\end{equation*}
	
	Condición ideal: Si $I_1 = I_2$, entonces $I_{min}=0$ y $\gamma = 1$ (contraste perfecto).
	La ecuación de interferencia puede reescribirse usando la visibilidad y la intensidad media local $I_T = I_1+I_2$:
	\begin{equation*}
		I = I_T [1 + \gamma \cos(\Delta\theta)]
	\end{equation*}
	
	\section{Coherencia}
	La coherencia mide la correlación y sincronización de las fases entre ondas. Para que dos ondas formen franjas estables, deben ser coherentes (relación de fase constante).
	\textit{[Imagen: comparacion\_luz\_coherente\_y\_no\_coherente.png]}
	
	Una onda perfectamente monocromática (una sola frecuencia) tendría una onda senoidal infinita, lo cual es una idealización matemática. En la realidad, las ondas se emiten en "trenes de onda" finitos.
	\textit{[Imagen: grafica\_tren\_ondas\_coherencia\_finita.png]}
	
	\subsection{Coherencia Temporal y Espacial}
	Existen dos dominios de la coherencia:
	
	\begin{center}
		\renewcommand{\arraystretch}{1.5}
		\begin{tabular}{|p{0.45\textwidth}|p{0.45\textwidth}|}
			\hline
			\textbf{\textcolor{blue}{COHERENCIA TEMPORAL (Longitudinal)}} & \textbf{\textcolor{blue}{COHERENCIA ESPACIAL (Transversal)}} \\
			\hline
			Correlación de fase del \textbf{mismo haz} en el \textbf{mismo punto} espacial, medido en dos \textbf{tiempos diferentes}. & Correlación de fase entre \textbf{dos puntos diferentes} del frente de onda, medido en el \textbf{mismo instante}. \\
			Relacionada con el ancho de banda del espectro de emisión (qué tan monocromática es la luz). & Relacionada con el tamaño físico de la fuente de luz. \\
			Se evalúa con un Interferómetro de Michelson. & Se evalúa con el experimento de doble rendija de Young. \\
			\hline
		\end{tabular}
	\end{center}
	\textit{[Imagen: comparativa\_ondas\_coherentes\_incoherentes\_espacial\_temporal.png]}
	
	\subsection{Cuantificación de la Coherencia Temporal}
	\begin{itemize}
		\item \textbf{Tiempo de coherencia ($\tau_c$):} Intervalo promedio en el cual la fase de la onda es predecible y mantiene su forma senoidal.
		\item \textbf{Longitud de coherencia ($L_c$):} Distancia física que recorre la luz durante su tiempo de coherencia. $L_c = c \tau_c$.
	\end{itemize}
	\textit{[Imagen: esquema\_tiempo\_de\_coherencia\_temporal.png]}
	
	Si un láser tiene un ancho de banda de frecuencias $\Delta \nu$, el tiempo de coherencia es $\tau_c = \frac{1}{\Delta \nu}$.
	Derivando de $\nu = c/\lambda$, sabemos que el ancho espectral en longitud de onda es $|\Delta \nu| = \frac{c}{\lambda^2} \Delta \lambda$.
	
	Sustituyendo esto, la longitud de coherencia queda relacionada con el ancho espectral ($\Delta \lambda$) y la pureza espectral ($Q = \lambda / \Delta \lambda$):
	\begin{equation*}
		L_c = \frac{c}{\Delta \nu} = \frac{\lambda^2}{\Delta \lambda} = \lambda Q
	\end{equation*}
	\textit{[Imagen: grafica\_ancho\_espectral\_delta\_lambda.png]}
	
	\section{Speckle (Moteado)}
	El "speckle" es un patrón de intensidad granular y aleatorio que aparece cuando luz coherente (láser) ilumina una superficie rugosa, cuyas variaciones de altura son del orden de la longitud de onda. 
	
	Cada punto de la superficie rugosa difracta la luz actuando como un emisor secundario. Al llegar al detector u ojo, todos estos haces secundarios con fases aleatorias interfieren.
	\textit{[Imagen: formacion\_speckle\_laser.png]}
	
	Estadísticamente, la probabilidad de encontrar una intensidad $I$ obedece a una distribución exponencial negativa:
	\begin{equation*}
		P(I) = \frac{1}{\langle I \rangle} e^{-I / \langle I \rangle}
	\end{equation*}
	Esto indica que el valor más probable de intensidad es cero (hay más motas oscuras que brillantes).
	\textit{[Imagen: histograma\_pdf\_intensidad\_speckle.png]}
	
	\subsection{Tipos de Speckle}
	\begin{itemize}
		\item \textbf{Speckle Objetivo:} Se forma directamente en el espacio libre de dispersión, al colocar un sensor sin lentes de por medio. El tamaño promedio del grano depende de la geometría óptica: $s = \frac{\lambda z}{D}$ (donde $D$ es el diámetro iluminado y $z$ la distancia a la pantalla).
		\item \textbf{Speckle Subjetivo:} Se forma al utilizar un sistema formador de imágenes (lentes u ojo). El tamaño del speckle está regido por el límite de difracción (disco de Airy) de la apertura de la lente: $\approx 2.4 \lambda u / D_{lente}$.
	\end{itemize}
	
	\section{Técnicas de Metrología y Pruebas No Destructivas (PND)}
	Las pruebas ópticas no destructivas permiten inspeccionar deformaciones, vibraciones, esfuerzos y defectos estructurales sin contacto físico y sin dañar la pieza. Se dividen en métodos coherentes (Interferometría, Holografía, Speckle) e incoherentes (Moiré, Proyección de franjas, Correlación digital DIC).
	
	\subsection{Interferometría de Speckle (SI) / ESPI}
	Originalmente visto como ruido, el speckle se aprovechó como portador de información. La interferometría de speckle compara la fase de un patrón de speckle \textbf{antes y después} de una deformación.
	\textit{[Imagen: diagrama\_bloques\_espi.png]}
	
	El patrón inicial tiene intensidad $I_A$ y el patrón tras aplicar carga mecánica es $I_B$:
	\begin{equation*}
		\Delta I = |I_A - I_B| = 4 a_R a_o \left| \sin\left(\varphi_o - \varphi_R + \frac{\Delta\varphi}{2}\right) \sin\left(\frac{\Delta\varphi}{2}\right) \right|
	\end{equation*}
	
	La sustracción digital de estas imágenes revela franjas de correlación de baja frecuencia asociadas estrictamente al desplazamiento del material $\Delta\varphi$.
	
	\subsubsection*{Vectores de Sensibilidad en ESPI}
	El cambio de fase $\Delta\phi$ detectado depende de la dirección de iluminación ($\vec{n}_i$) y observación ($\vec{n}_o$) respecto al vector de desplazamiento real del objeto ($\vec{d}$):
	\begin{equation*}
		\Delta\phi = (\vec{K}_i - \vec{K}_o)\cdot\vec{d} = \frac{2\pi}{\lambda}(\vec{n}_i - \vec{n}_o)\cdot\vec{d}
	\end{equation*}
	
	\textbf{1. Configuración Fuera del Plano (Out-of-plane):} Luz colineal con la observación. Sensible al desplazamiento $z$.
	$\Delta\phi = \frac{4\pi}{\lambda}d_z$
	\textit{[Imagen: esquema\_optico\_dspi\_fuera\_plano.png]}
	
	\textbf{2. Configuración En el Plano (In-plane):} Dos haces simétricos iluminan con ángulo $\alpha$.
	$\Delta\phi = \frac{4\pi}{\lambda} \sin(\alpha) d_x$
	\textit{[Imagen: esquema\_optico\_dspi\_en\_plano.png]}
	
	\textbf{3. Cizalladura (Shearography):} Se usa un prisma divisor; interfieren dos puntos contiguos de la misma superficie ($\Delta x, \Delta y$). Mide directamente el gradiente de deformación, ideal para encontrar defectos internos.
	$\Delta\phi \approx \frac{4\pi}{\lambda} \left[ \frac{\partial d_z}{\partial x} \right] \Delta x$
	\textit{[Imagen: esquema\_optico\_shearography\_prisma.png]}
	
	\section{Interferometría de Medición de Fase (PMI y PSI)}
	La simple observación de franjas de interferencia presenta ambigüedad de fase (no se distingue si el orden de franja sube o baja ya que $\cos(\phi) = \cos(-\phi)$). Las técnicas de demodulación resuelven esto.
	
	\subsection{Métodos Temporales (Phase Shifting / Phase Stepping)}
	Se toman secuencialmente múltiples interferogramas, introduciendo un retraso de fase calibrado y controlado (ej. usando un espejo en un actuador piezoeléctrico - PZT).
	\begin{equation*}
		I_n(x,y) = I_{dc} + I_{ac} \cos[\varphi(x,y) + \alpha_n]
	\end{equation*}
	
	\textbf{Algoritmo de 4 pasos} ($\alpha_n = 0, \pi/2, \pi, 3\pi/2$):
	Resolviendo el sistema de ecuaciones para las cuatro intensidades $I_1, I_2, I_3, I_4$:
	\begin{equation*}
		\varphi(x,y) = \arctan \left[ \frac{I_4(x,y) - I_2(x,y)}{I_1(x,y) - I_3(x,y)} \right]
	\end{equation*}
	El cálculo proporciona el mapa de fase módulo $2\pi$ (fase envuelta). A partir de ahí, la deformación real $d(x,y)$ se calcula como: $d(x,y) = \frac{\lambda}{4\pi} \varphi(x,y)$.
	
	\subsection{Métodos Espaciales (Transformada de Fourier / Takeda)}
	En lugar de tomar múltiples imágenes, se capta un único interferograma utilizando una configuración "fuera de eje" (off-axis) que inyecta una alta frecuencia portadora de franjas muy finas.
	\begin{equation*}
		I(x,y) = a(x,y) + b(x,y) \cos[\phi(x,y) + 2\pi f_0 x]
	\end{equation*}
	
	Al aplicar la Transformada Rápida de Fourier (FFT), el espectro de frecuencias muestra un lóbulo central (fondo) y dos lóbulos laterales (fase modulada).
	\begin{equation*}
		I(u,v) = A(u,v) + C(u,v) + C^*(-u,-v)
	\end{equation*}
	Se filtra espacialmente (aislando el lóbulo $C(u,v)$), se centra y se aplica la Transformada Inversa para calcular directamente el mapa de fase:
	\begin{equation*}
		\phi(x, y) = \arctan \left[ \frac{\text{Im}(c(x,y))}{\text{Re}(c(x,y))} \right]
	\end{equation*}
	
	\section{Ejercicios Propuestos y Anotaciones Adicionales}
	
	\subsection{Ejercicios de Índices de Refracción}
	El índice de refracción de la glicerina y el diamante con respecto al aire son 1.46 y 2.42 respectivamente. Calcula la velocidad y la longitud de onda de la luz en la glicerina y en el diamante para luz que viaja a través del aire con $\lambda = 5400 \text{ \AA}$. Posteriormente, calcular el índice de refracción relativo del diamante con respecto a la glicerina.
	
	\subsection{Ejercicios de Superposición y Fase}
	Dada la onda $E(t) = E_0\cos(\omega t + \phi)$, encuentre la velocidad de oscilación $V$, el valor de fase y la amplitud inicial, evaluando $E(0)$ y $V(0)$. 
	
	\textit{Resolución propuesta de las notas:}
	\begin{align*}
		V(t) &= \frac{dE}{dt} = - \omega E_0 \sin(\omega t + \phi) \\
		\text{Evaluando en } t=0: \quad E(0) &= E_0 \cos(\phi) \quad \text{y} \quad V(0) = - \omega E_0 \sin(\phi) \\
		\frac{V(0)}{E(0)} &= \frac{- \omega E_0 \sin(\phi)}{E_0 \cos(\phi)} = - \omega \tan(\phi) \implies \phi = \arctan\left( \frac{V(0)}{-\omega E(0)} \right)
	\end{align*}
	Para recuperar la amplitud $E_0$:
	\begin{equation*}
		E^2(0) + \frac{V^2(0)}{\omega^2} = E_0^2 (\cos^2\phi + \sin^2\phi) = E_0^2
	\end{equation*}
	
	\subsection{Ejercicios de Interferencia y Caminos Ópticos}
	Dos haces de microondas de longitud de onda $6.00 \times 10^{-3}$ m recorren 5 cm hasta un detector. Uno pasa por aire y el otro por cuarzo ($n = 1.54$). ¿Provocan interferencia constructiva o destructiva?
	
	\textit{Resolución:}
	\begin{align*}
		\text{OPD} &= (n_{\text{cuarzo}} - n_{\text{aire}}) \times z \\
		&= (1.54 - 1.00) \times 0.05 \text{ m} = 0.027 \text{ m}
	\end{align*}
	Dividiendo el OPD entre la longitud de onda:
	\begin{equation*}
		\text{Órdenes de interferencia} = \frac{0.027}{6.00 \times 10^{-3}} = 4.5
	\end{equation*}
	Como resulta en un número entero más medio ciclo ($4 + 0.5$), las ondas llegarán en contrafase, causando interferencia \textbf{destructiva}.
	
	\subsection{Ejercicios de Coherencia}
	\begin{enumerate}
		\item Determine el ancho de banda ($\Delta \nu$), longitud de coherencia ($L_c$) y tiempo de coherencia ($\tau_c$) de la luz blanca asumiendo su espectro entre 384 THz y 769 THz. \textit{(Desarrollado en las notas previas: $L_c \approx 779$ nm).}
		\item La luz de $\lambda = 4800 \text{ \AA}$ tiene una longitud de coherencia equivalente a 25 ondas. ¿Cuál es la longitud de coherencia absoluta y el tiempo de coherencia?
		\item Se modula un láser continuo ($\lambda = 6500 \text{ \AA}$) en pulsos de 0.5 ns usando un obturador. Calcule $L_c$, $\Delta \nu$, y $\Delta \lambda$.
	\end{enumerate}
	
\end{document}